\section{Két munkahengeres vezérlés. Két munkahenger alternáló mozgása, ahol a munkahengerek mozgása egymáshoz képest egy lépéssel el van tolva. (Út-lépés diagram) (Indítás és leállítás 3/2-es szelepekkel és a munkahengereknél végálláskapcsolókkal)}

\begin{figure}[H]
    \centering
    \begin{tikzpicture}[thick, >=Stealth, scale=0.7, transform shape]
        \begin{scope}[shift={(1,0)}]
        % Cylinders
        \begin{scope}[shift={(-1, 8)}]
            \cyDouble
            \draw (2.8, 1.2) -- (2.8, 2.2) node[above] {1S1};
            \draw (5.0, 1.2) -- (5.0, 2.2) node[above] {1S2};
        \end{scope}
        \begin{scope}[shift={(5, 8)}]
            \cyDouble
            \draw (2.8, 1.2) -- (2.8, 2.2) node[above] {2S1};
            \draw (5.0, 1.2) -- (5.0, 2.2) node[above] {2S2};
        \end{scope}
            
            \begin{scope}[shift={(-2.5, 1.5)}]
                \valveFiveTwo
                \actPilotLeft
                \actPilotRight
                \portExhaust{1.3}{0}
                \portSupply{1.5}{0}
                \portExhaust{1.7}{0}
                \node[below] at (1, -0.8) {5/2 (A)};
                \draw (1.3, 1) |- (0.5, 3) -- (0.5, 6.0) ;
                \draw (1.7, 1) |- (3.0, 4) -- (3.0, 6.0) ;
            \end{scope}
            
            \begin{scope}[shift={(4.5, 1.5)}]
                \valveFiveTwo
                \actPilotLeft
                \actPilotRight
                \portExhaust{1.3}{0}
                \portSupply{1.5}{0}
                \portExhaust{1.7}{0}
                \node[below] at (1, -0.8) {5/2 (B)};
                \draw (1.3, 1) |- (-0.5, 3) -- (-0.5, 6.0) ;
                \draw (1.7, 1) |- (2.0, 4) -- (2.0, 6.0) ;
            \end{scope}
            
            \draw[dashed] (-3, -2) rectangle (8, 0);
            \node at (2.5, -1) {Lejátszó logika \& 3/2 Görgős szelepek lánca};

            \draw (-2.5, 0) -- (-2.5, 1.5);
            \draw (-0.5, 0) -- (-0.5, 1.5);
            \draw (4.5, 0) -- (4.5, 1.5);
            \draw (6.5, 0) -- (6.5, 1.5);
        \end{scope}
    \end{tikzpicture}
\end{figure}
